%----------------------------------------------------------------------------------------
%	PACKAGES AND THEMES
%----------------------------------------------------------------------------------------
\documentclass[aspectratio=169,xcolor=dvipsnames, t]{beamer}
\usepackage[utf8]{inputenc}
\usepackage[vietnamese]{babel}
\usepackage{ragged2e}
\usepackage{fontspec} % Allows using custom font. MUST be before loading the theme!
\usetheme{SimplePlusAIC}
\usepackage{hyperref}
\usepackage{graphicx} % Allows including images
\usepackage{booktabs} % Allows the use of \toprule, \midrule and  \bottomrule in tables
\usepackage{svg} %allows using svg figures
\usepackage{tikz}
\usepackage{makecell}
\usepackage{wrapfig}
% ADD YOUR PACKAGES BELOW

%----------------------------------------------------------------------------------------
%	TITLE PAGE CONFIGURATION
%----------------------------------------------------------------------------------------

\title[Hệ thống gợi ý việc làm]{PHÁT TRIỂN HỆ THỐNG TUYỂN DỤNG TRỰC TUYẾN TÍCH HỢP CHỨC NĂNG GỢI Ý VIỆC LÀM SỬ DỤNG MÔ HÌNH HỌC SÂU} % The short title appears at the bottom of every slide, the full title is only on the title page
\subtitle{ĐỒ ÁN TỐT NGHIỆP ĐẠI HỌC}
\author{%
Giáo viên hướng dẫn: PGS.TS. Trần Đình Quế\\[0.3em]
Sinh viên thực hiện: Đào Duy Thông\\[0.1em]
\hspace*{3.45cm} Nguyễn Ngọc Hà\\[0.1em]
\hspace*{3.45cm} Nguyễn Minh Tùng%
}
\institute[]{Khoa CNTT 1 - PTIT
}
% Your institution as it will appear on the bottom of every slide, maybe shorthand to save space


\date{2025} % Date, can be changed to a custom date
%----------------------------------------------------------------------------------------
%	PRESENTATION SLIDES
%----------------------------------------------------------------------------------------
\begin{document}

\maketitlepage

\begin{frame}[t]{Nội dung}
    % Throughout your presentation, if you choose to use \section{} and \subsection{} commands, these will automatically be printed on this slide as an overview of your presentation
    \tableofcontents
\end{frame}

%------------------------------------------------
% Section divider frame
\makesection{Mở đầu}
%------------------------------------------------

\begin{frame}{Bối cảnh xã hội}
\justifying

\textbf{Chuyển đổi số trong tuyển dụng:}
\begin{itemize}
    \item Các nền tảng như VietnamWorks, TopCV, CareerBuilder tạo hệ sinh thái kết nối rộng lớn.
    \item Quy mô thị trường tuyển dụng trực tuyến Việt Nam: \textbf{~1,2 tỷ USD}.
\end{itemize}

\vspace{0.3cm}
\textbf{Thách thức hiện tại:}
\begin{itemize}
    \item Quá tải thông tin: hàng nghìn tin tuyển dụng mỗi ngày, vô số hồ sơ ứng viên.
    \item Khó khăn trong việc tìm đúng người cho đúng vị trí.
    \item Các hệ thống truyền thống chưa cá nhân hóa hiệu quả $\rightarrow$ bỏ lỡ cơ hội.
\end{itemize}

\vspace{0.3cm}
\textbf{Cơ hội:} AI và Deep Learning giúp hiểu ngữ nghĩa, gợi ý chính xác hơn.

\end{frame}

\begin{frame}{Nội dung nghiên cứu}
\justifying

\begin{enumerate}
    \item \textbf{Khảo sát và phân tích hệ thống tuyển dụng hiện có:} Khảo sát các chức năng, nghiệp vụ, phân tích ưu điểm, hạn chế và xu hướng ứng dụng trí tuệ nhân tạo trong bài toán gợi ý việc làm.
    \item \textbf{Phân tích yêu cầu và thiết kế hệ thống:} Thực hiện phân tích yêu cầu và thiết kế hệ thống thông qua các mô hình UML như: biểu đồ ca sử dụng, biểu đồ lớp, biểu đồ tuần tự, …; thiết kế kiến trúc tổng thể của hệ thống, thiết kế chi tiết API, cơ sở dữ liệu, giao diện người dùng.
\end{enumerate}

\end{frame}
\begin{frame}{Nội dung nghiên cứu}
\justifying

\begin{enumerate}
    \setcounter{enumi}{2}
    \item \textbf{Phát triển mô đun gợi ý việc làm :} Nghiên cứu các mô hình học sâu phù hợp cho bài toán gợi ý việc làm. Ứng dụng mô hình học sâu để xây dựng module gợi ý việc làm, đánh giá hiệu quả gợi ý.
    \item \textbf{Phát triển và đánh giá hệ thống:} Phát triển hệ thống web hoàn chỉnh bao gồm giao diện người dùng, API, cơ sở dữ liệu và mô-đun gợi ý. Triển khai hệ thống lên môi trường đám mây.
\end{enumerate}

\end{frame}
% CHƯƠNG 2: PHÂN TÍCH VÀ THIẾT KẾ HỆ THỐNG

\makesection{Phân tích và Thiết kế hệ thống}

% 1. YÊU CẦU CHỨC NĂNG (trang 13 báo cáo, 3 gạch đầu dòng)

\begin{frame}{Yêu cầu chức năng}
\justifying

Hệ thống phục vụ ba nhóm người dùng chính:

\vspace{0.3cm}
\begin{itemize}
    \item \textbf{Ứng viên:} Đăng ký/đăng nhập, tạo và quản lý CV, tìm kiếm việc làm/công ty, lưu việc làm, xem việc làm được gợi ý, nộp hồ sơ trực tuyến, theo dõi đơn ứng tuyển.
    
    \item \textbf{Nhà tuyển dụng:} Quản lý thông tin công ty, quản lý thành viên và vai trò trong công ty, tạo/chỉnh sửa/đóng tin tuyển dụng, tiếp nhận và đánh giá đơn ứng tuyển, tìm kiếm và lưu hồ sơ ứng viên phù hợp.
    
    \item \textbf{Quản trị viên:} Quản lý tài khoản, phân quyền, kiểm duyệt thông tin công ty, quản lý mẫu CV, đảm bảo tính minh bạch và tuân thủ quy định sử dụng nền tảng.
\end{itemize}

\end{frame}

% 2. YÊU CẦU PHI CHỨC NĂNG (check lại)

\begin{frame}{Yêu cầu phi chức năng}
\justifying

\begin{itemize}
    \item \textbf{Hiệu năng}: Thời gian phản hồi trung bình cho các thao tác cơ bản không vượt quá 3 giây.
    
    \item \textbf{Bảo mật}: 
    \begin{itemize}
        \item Bảo vệ dữ liệu cá nhân của người dùng và thông tin
doanh nghiệp
        \item Áp dụng cơ chế xác thực an toàn và mã hoá dữ liệu nhạy
cảm
        \item Phân quyền truy cập rõ ràng cho từng loại người dùng
    \end{itemize}
    
    \item \textbf{Khả năng mở rộng }: 
    \begin{itemize}
        \item Cho phép mở rộng hệ thống dễ dàng khi số lượng người dùng tăng
        \item Hỗ trợ tích hợp thêm các tính năng mới mà không ảnh
hưởng tới hệ thống hiện tại
    \end{itemize}
    
    \item \textbf{Khả năng sử dụng }: 
    \begin{itemize}
        \item Giao diện trực quan, thân thiện, dễ thao tác với cả người dùng phổ thông
        \item Hỗ trợ hiển thị tốt trên cả máy tính và thiết bị di động
    \end{itemize}
    
\end{itemize}

\end{frame}

% 3. MÔ HÌNH USECASE (giữ nguyên)

\begin{frame}{Mô hình chức năng - Usecase}
\begin{figure}
    \centering
    \includegraphics[width=0.75\textwidth]{figures/Overview.png}
    \caption{Usecase tổng quan}
\end{figure}
\end{frame}

% 4. ACTIVITY DIAGRAM (2 hình 1 trang) - THÊM MỚI

\begin{frame}{Mô hình hoạt động - Activity Diagram}
\justifying
\small

\begin{columns}[T]
\column{0.45\textwidth}

\textbf{Activity Diagram} mô tả các quy trình, luồng công việc và hoạt động trong hệ thống:

\vspace{0.2cm}
\begin{itemize}
    \item Hiển thị các bước tuần tự mà người dùng và hệ thống thực hiện
    \item Xác định điểm quyết định, nhánh điều kiện
    \item Giúp xác nhận rõ ràng logic nghiệp vụ
\end{itemize}



\column{0.52\textwidth}
\begin{figure}
    \centering
    \includegraphics[width=0.5\textwidth]{figures/CandidateCreateCV1.png}
    \caption{Biểu đồ hoạt động: Ứng viên tạo CV}
\end{figure}
\end{columns}
\end{frame}

% \begin{frame}{Mô hình hoạt động (tiếp theo)}
% \begin{figure}
%     \centering
%     \includegraphics[width=0.3\textwidth]{figures/CandidateApplyJob1.png}
%     \caption{Biểu đồ hoạt động: Ứng viên ứng tuyển việc làm}
% \end{figure}

% \end{frame}
% 5. CLASS DIAGRAM - Thêm giải thích trước (tr 28) - ĐẢO VỊ TRÍ

\begin{frame}{Xác định lớp, thuộc tính và hoạt động}
\justifying
\small

Dựa trên phân tích yêu cầu nghiệp vụ và các usecase, xác định được các lớp thực thể chính:

\vspace{0.3cm}
\textbf{5 lớp cốt lõi:}
\begin{itemize}
    \item \textbf{User}: Người dùng (ứng viên, nhà tuyển dụng, admin) với các thuộc tính: id, email, passwordHash, fullName, role, status...
    \item \textbf{Company}: Công ty với thông tin như name, website, description, industry, companySize, status...
    \item \textbf{CV}: Hồ sơ ứng viên với title, fullName, email, summary, titleEmbedding (vector AI)...
    \item \textbf{Job}: Tin tuyển dụng với title, description, location, experienceLevel, titleEmbedding...
    \item \textbf{Application}: Đơn ứng tuyển kết nối CV và Job, với status, coverLetter, notes...
\end{itemize}

\vspace{0.2cm}
\textbf{Các lớp phụ trợ:} CVSkill, Education, WorkExperience, Salary, JobBenefit, CompanyMember, 

SavedJob, SavedCV, Notification...

\end{frame}

\begin{frame}{Mô hình cấu trúc - Class Diagram}
\begin{figure}
    \centering
    \includegraphics[width=0.65\textwidth]{figures/class.png}
    \caption{Biểu đồ lớp}
\end{figure}
\end{frame}

% 6. SEQUENCE DIAGRAM (2 hình) - THÊM MỚI

\begin{frame}{Mô hình hành vi - Sequence Diagram}
\justifying
\small

\begin{columns}[T]
\column{0.40\textwidth}

\textbf{Sequence Diagram} mô tả tương tác giữa các đối tượng theo thời gian:

\vspace{0.2cm}
\begin{itemize}
    \item Hiển thị thứ tự gọi giữa các thành phần (Actor, Controller, Service, Database)
    \item Xác định luồng dữ liệu và phản hồi
    \item Giúp hiểu rõ cách các module phối hợp
\end{itemize}


\column{0.55\textwidth}
\begin{figure}
    \centering
    \includegraphics[width=0.63\textwidth]{figures/CandidateCreateCV.png}
    \caption{Biểu đồ tuần tự: Ứng viên tạo CV}
\end{figure}
\end{columns}
\end{frame}

% \begin{frame}{Mô hình hành vi - Sequence Diagram (tiếp theo)}
% \begin{figure}
%     \centering
%     \includegraphics[width=0.6\textwidth]{figures/CandidateApplyJob.png}
%     \caption{Biểu đồ tuần tự: Ứng viên ứng tuyển việc làm}
% \end{figure}
% \end{frame}

% 7. STATE DIAGRAM (2 hình) - THÊM MỚI

\begin{frame}{Mô hình trạng thái - State Machine Diagram}
\justifying
\small

\begin{columns}[T]
\column{0.4\textwidth}

\textbf{State Machine Diagram} mô tả các trạng thái và chuyển đổi giữa chúng:

\vspace{0.2cm}
\begin{itemize}
    \item Hiển thị tất cả trạng thái có thể của một thực thể
    \item Xác định điều kiện kích hoạt (trigger) chuyển đổi trạng thái
    \item Giúp xác thực logic trạng thái trong quy trình
\end{itemize}

\column{0.55\textwidth}
\begin{figure}
    \centering
    \includegraphics[width=0.75\textwidth]{figures/Application.png}
    \caption{Biểu đồ trạng thái: Đơn ứng tuyển}
\end{figure}
\end{columns}
\end{frame}

% \begin{frame}{Mô hình trạng thái (tiếp theo)}
% \begin{figure}
%     \centering
%     \includegraphics[width=0.5\textwidth]{figures/Company.png}
%     \caption{Biểu đồ trạng thái: Công ty}
% \end{figure}
% \end{frame}

% 8. CLEAN ARCHITECTURE (2 trang: 1 lý thuyết + 1 hình) - THÊM MỚI

\begin{frame}{Clean Architecture - Tổng quan}
\begin{columns}[T]
\column{0.55\textwidth}
\justifying

\textbf{Clean Architecture} là mô hình kiến trúc phần mềm tách biệt các tầng nghiệp vụ, giúp hệ thống độc lập với framework, database và UI.

\vspace{0.3cm}
\textbf{Mô hình 4 tầng:}
\begin{enumerate}
    \item \textbf{Domain Layer}: Nghiệp vụ cốt lõi, hoàn toàn độc lập
    \item \textbf{Application Layer}: Logic nghiệp vụ cụ thể
    \item \textbf{Interface Adapters}: Chuyển đổi dữ liệu giữa tầng nghiệp vụ và bên ngoài
    \item \textbf{Infrastructure}: Triển khai cụ thể (DB, API, Services)
\end{enumerate}

\column{0.4\textwidth}
\begin{figure}
    \centering
    \includegraphics[width=\textwidth]{figures/clean-architecture.jpeg}
    \caption{Kiến trúc Clean Architecture}
\end{figure}
\end{columns}
\end{frame}

\begin{frame}{Thiết kế hệ thống theo Clean Architecture}
\begin{columns}[T]
\column{0.45\textwidth}
\justifying

\vspace{0.2cm}
\textbf{1. Domain Layer:} Entities: User, Job, Company,...

\vspace{0.2cm}
\textbf{2. Application Layer:} Use Cases: CreateCV, ApplyJob,...

\vspace{0.2cm}
\textbf{3. Interface Adapters:}
\begin{itemize}
    \item Controllers: REST API
    \item Repositories: Data access
    \item Presenters: DTO transformation
\end{itemize}

\vspace{0.2cm}
\textbf{4. Infrastructure:} PostgreSQL, Firebase, AI Service, File Storage

\column{0.5\textwidth}
\begin{figure}
    \centering
    \includegraphics[width=\textwidth]{figures/Layer.png}
    \caption{Kiến trúc phân tầng của hệ thống}
\end{figure}
\end{columns}
\end{frame}

% 9. CONTRACT + METHOD: createCV - THÊM MỚI

\begin{frame}[fragile]{Hợp đồng và phương thức: createCV}
\justifying
\small

\begin{columns}[T]
\column{0.48\textwidth}
\textbf{Use Case}: CreateCVUseCase

\vspace{0.2cm}
\textbf{Hợp đồng (Contract):}
\begin{itemize}
    \item \textbf{Tiền điều kiện}: User đã đăng nhập, có role = CANDIDATE
    \item \textbf{Hậu điều kiện}: CV mới được tạo và lưu vào database, trả về id của CV
\end{itemize}

\vspace{0.2cm}
\textbf{Chữ ký phương thức:}
\begin{verbatim}
createCV(userId: string, 
cvData: CreateCVDto): 
Promise<CV>
\end{verbatim}

\column{0.48\textwidth}
\textbf{Thuật toán:}
\begin{enumerate}
    \item Validate dữ liệu đầu vào (title, fullName, email, phoneNumber...)
    \item Kiểm tra userId có tồn tại và role = CANDIDATE
    \item Tạo đối tượng CV với dữ liệu từ cvData
    \item Lưu CV vào database qua CVRepository
    \item Gọi AI Service để tạo embedding cho title, skills, experience
    \item Cập nhật CV với embedding vectors
    \item Return CV object
\end{enumerate}
\end{columns}

\end{frame}

% 10. ERD (CẢ ORM) - ĐẢO VỊ TRÍ LÊN TRƯỚC DEPLOYMENT

\begin{frame}{Thiết kế cơ sở dữ liệu}
\begin{figure}
    \centering
    \includegraphics[width=0.65\textwidth]{figures/orm.png}
    \caption{ORM Diagram}
\end{figure}
\end{frame}

\begin{frame}{Thiết kế cơ sở dữ liệu}
\begin{figure}
    \centering
    \includegraphics[width=0.75\textwidth]{figures/erd.png}
    \caption{Entity Relationship Diagram}
\end{figure}

\vspace{-0.3cm}
\textbf{CSDL quan hệ}: PostgreSQL với 20+ bảng chính, bao gồm User, Company, CV, Job, Application, Education, WorkExperience, Salary, CompanyMember, SavedJob...

\end{frame}

% 11. GIAO DIỆN: Lý thuyết + 1 hình - THÊM MỚI

\begin{frame}{Thiết kế giao diện người dùng}
\begin{columns}[T]
\column{0.55\textwidth}
\justifying
\small

\textbf{Đối tượng người dùng:} Ứng viên (18-35 tuổi) và Nhà tuyển dụng (HR chuyên nghiệp)

\vspace{0.2cm}
\textbf{Cơ sở thiết kế:} ISO 9241-210 - lấy người dùng làm trung tâm

\vspace{0.2cm}
\textbf{Ba nguyên tắc chính:}
\begin{enumerate}
    \item \textbf{Mã hóa màu sắc}: Xanh dương (tin tưởng), xanh lá (tích cực), đỏ (cảnh báo)
    \item \textbf{Nhất quán điều hướng}: Ứng viên dùng thanh ngang, Nhà tuyển dụng dùng thanh dọc
    \item \textbf{Phân cấp thị giác}: Cỡ chữ, tương phản cao, thông tin theo khối
\end{enumerate}

\column{0.4\textwidth}
\begin{figure}
    \centering
    \includegraphics[width=0.75\textwidth]{figures/Candidate Home.png}
    \caption{Giao diện trang chủ ứng viên}
\end{figure}
\end{columns}
\end{frame}

% 12. DEPLOYMENT DIAGRAM - GIỮ NGUYÊN NHƯNG ĐẢO XUỐNG CUỐI

\begin{frame}{Kiến trúc vật lý - Deployment}
\begin{figure}
    \centering
    \includegraphics[width=0.54\textwidth]{figures/Deployment Diagram.png}
    \caption{Sơ đồ triển khai hệ thống}
\end{figure}
\end{frame}

% CHƯƠNG 3: KỸ THUẬT AI GỢI Ý VIỆC LÀM
% File này chứa nội dung đã được sửa đổi theo yêu cầu

\makesection{Kỹ thuật AI gợi ý việc làm}

% PHẦN 1: CƠ SỞ LÝ THUYẾT

%--- 1.1 BERT & PhoBERT - Nền tảng ---
\begin{frame}{BERT \& PhoBERT – Biểu diễn ngữ cảnh ở mức Token}
\justifying
\small

\textbf{BERT là gì?}
\begin{itemize}
    \item Mô hình Transformer Encoder học biểu diễn văn bản theo ngữ cảnh hai chiều
    \item Được \textbf{pre-train} (huấn luyện trước) bằng Masked Language Modeling:
    \begin{itemize}
        \item Che ngẫu nhiên 15\% từ trong câu, mô hình dự đoán từ bị che
        \item Học được ngữ cảnh từ cả hai phía (trái và phải)
    \end{itemize}
\end{itemize}

% \vspace{0.2cm}
% \textbf{Token-level embeddings:} BERT sinh \textit{1 vector riêng cho mỗi từ}
% \begin{itemize}
%     \item Input: \texttt{"Tôi là lập trình viên Python"}
%     \item Output: \texttt{[CLS]} → vec$_1$, \texttt{Tôi} → vec$_2$, \texttt{là} → vec$_3$, ..., \texttt{Python} → vec$_6$, \texttt{[SEP]} → vec$_7$
%     \item Mỗi vec có 768 chiều, phản ánh nghĩa của từ trong ngữ cảnh
% \end{itemize}

\vspace{0.2cm}
\textbf{BERT dùng cho tác vụ gì?}
\begin{itemize}
    \item Phân loại văn bản: dùng vector \texttt{[CLS]} → FC layer → nhãn
    \item NER: dùng vector mỗi token để gán nhãn (PERSON, ORG, LOC...)
    \item Question Answering: tìm token bắt đầu/kết thúc câu trả lời
\end{itemize}

\end{frame}

%--- 1.2 Sentence-BERT & SimCSE - Phương pháp ---
\begin{frame}{Từ Token-level đến Sentence-level Embeddings}
\justifying
\small

\textbf{Vấn đề khi dùng BERT cho sentence similarity:}
\begin{itemize}
    \item BERT output: \textit{N tokens} × 768 dim → cần gộp (mean/max pooling hoặc [CLS])
    \item \textbf{Nhưng BERT không được train} để vector sau pooling phản ánh similarity!
    \item Mean pooling thủ công → kém chính xác, không tối ưu
\end{itemize}

\vspace{0.3cm}
\textbf{Sentence-BERT:} Fine-tune BERT để học sentence embeddings
\begin{itemize}
    \item Kiến trúc Siamese: hai câu cùng encoder (shared weights)
    \item Train với classification/regression objective → embeddings có ý nghĩa similarity
    \item Encode 1 lần, so sánh nhanh bằng cosine similarity
\end{itemize}

\vspace{0.3cm}
\textbf{SimCSE – Contrastive Learning đơn giản:}
\begin{itemize}
    \item \textbf{Supervised:} Dùng NLI data (entailment = positive, contradiction = negative)
    \item Loss: Kéo câu cùng nghĩa gần, đẩy câu khác nghĩa xa
\end{itemize}

\end{frame}

%--- 1.3 Mô hình VoVanPhuc - Áp dụng ---
\begin{frame}{Mô hình VoVanPhuc/sup-SimCSE-PhoBERT}
\justifying
\small

\textbf{Kiến trúc:} PhoBERT + Supervised SimCSE
\begin{itemize}
    \item \textbf{Base model:} PhoBERT-base (110M parameters)
    \item \textbf{Training:} Supervised SimCSE trên Vietnamese NLI dataset
    \item \textbf{Output:} Sentence embeddings 768 chiều, chuẩn hóa L2
\end{itemize}

\vspace{0.3cm}
\textbf{Ưu điểm cho hệ thống gợi ý:}
\begin{itemize}
    \item Tối ưu cho tiếng Việt (PhoBERT backbone)
    \item Embedding ở mức câu → phù hợp cho matching CV–Job
    \item Cosine similarity phản ánh độ tương đồng ngữ nghĩa tốt
    \item Inference nhanh: encode 1 lần, tái sử dụng nhiều lần
\end{itemize}


\end{frame}



%--- 1.2 FAISS ---
\begin{frame}{Cơ sở lý thuyết: FAISS}
\justifying
\small

\textbf{FAISS} (Facebook AI Similarity Search) - Thư viện tìm kiếm vector hiệu quả:

\vspace{0.3cm}
\textbf{Các loại Index:}
\begin{itemize}
    \item \textbf{Flat Index}: Tìm kiếm vét cạn, chính xác 100\% nhưng O(n).
    \item \textbf{IVF (Inverted File)}: Phân cụm vector, chỉ tìm trong cụm gần nhất $\rightarrow$ O(log n).
    \item \textbf{HNSW}: Đồ thị đa tầng, nhanh nhất nhưng tốn RAM.
\end{itemize}

\textbf{Hệ thống sử dụng IVF} vì cân bằng về tốc độ, tài nguyên sử dụng và hiệu năng tìm kiếm.

\vspace{0.3cm}
\textbf{Độ đo tương đồng Cosine:}
$$\text{cosine}(\mathbf{a}, \mathbf{b}) = \frac{\mathbf{a} \cdot \mathbf{b}}{||\mathbf{a}|| \cdot ||\mathbf{b}||} \in [-1, 1]$$

\end{frame}

%--- 1.3 LỌC PHÂN TẦNG ---
\begin{frame}{Cơ sở lý thuyết: Lọc phân tầng}
\justifying
\small

\textbf{Cascade Filtering}: Xử lý dữ liệu theo nhiều giai đoạn tuần tự, từ thô đến tinh.

\vspace{0.3cm}
\textbf{Nguyên lý:}
\begin{itemize}
    \item Giai đoạn đầu: Tiêu chí đơn giản, chi phí thấp, loại bỏ nhanh phần lớn ứng viên không phù hợp.
    \item Giai đoạn sau: Tiêu chí phức tạp hơn, xử lý trên tập nhỏ đã được lọc.
    \item Kết quả: Tối ưu chi phí tính toán tổng thể.
\end{itemize}

% \vspace{0.3cm}
% \textbf{Ứng dụng trong hệ thống:}
% \begin{enumerate}
%     \item \textbf{Vòng 1 - Tiêu đề}: FAISS tìm kiếm K1 jobs theo title embedding.
%     \item \textbf{Vòng 2 - Kinh nghiệm}: So khớp experience với requirements $\rightarrow$ K2 jobs.
%     \item \textbf{Vòng 3 - Kỹ năng}: So khớp skills embedding $\rightarrow$ K jobs cuối cùng.
% \end{enumerate}

\begin{figure}
    \centering
    \includegraphics[width=0.9\textwidth]{figures/Lọc phân tầng.png}
    \caption{Kiến trúc lọc phân tầng ba giai đoạn}
\end{figure}

\end{frame}

% PHẦN 2: CÀI ĐẶT

%--- 2.1 KIẾN TRÚC TỔNG THỂ ---
\begin{frame}{Cài đặt: Kiến trúc tổng thể}
Hệ thống gợi ý việc làm được thiết kế theo kiến trúc module hóa, trong đó các thành phần được tổ chức thành các module độc lập với trách nhiệm rõ ràng.
\begin{figure}
    \centering
    \includegraphics[width=0.65\textwidth]{figures/AI-Architecture.png}
    \caption{Kiến trúc hệ thống AI gợi ý}
\end{figure}

% \vspace{-0.3cm}
% \textbf{Các thành phần:}
% \begin{itemize}
%     \item \textbf{Embedding Service}: PhoBERT-SimCSE tạo vector 768 chiều.
%     \item \textbf{Vector Store}: PostgreSQL (persistent) + FAISS (in-memory, fast).
%     \item \textbf{Recommendation Engine}: Cascade filtering 3 vòng.
%     \item \textbf{API}: Flask RESTful API với endpoint \texttt{POST /recommend}.
% \end{itemize}

\end{frame}

%--- 2.2 CHIẾN LƯỢC NHÚNG VĂN BẢN ---
\begin{frame}{Cài đặt: Chiến lược nhúng văn bản}
\justifying
\small

\textbf{Nhúng riêng từng trường:} CV và Job được nhúng thành 3 vectors riêng biệt

\vspace{0.2cm}
\textbf{Lợi ích:}
\begin{itemize}
    \item So sánh có mục đích: title-title, skills-skills
    \item Kiểm soát trọng số từng yếu tố trong lọc phân tầng
    \item Linh hoạt khi dữ liệu thiếu
\end{itemize}

\begin{figure}
    \centering
    \includegraphics[width=0.55\textwidth]{figures/Chiến lược nhúng.png}
    \caption{Các trường được nhúng và ánh xạ trong lọc phân tầng}
\end{figure}

\end{frame}

% %--- 2.3 LƯU TRỮ ---
% \begin{frame}{Cài đặt: Lưu trữ vector}
% \justifying

% \textbf{Chiến lược lưu trữ kép:}

% \vspace{0.3cm}
% \textbf{1. PostgreSQL (Persistent Storage):}
% \begin{itemize}
%     \item Lưu vector dưới dạng JSONB/Bytea
%     \item Bảng CV: (id, titleEmbedding, skillsEmbedding, experienceEmbedding)
%     \item Bảng Job: (id, titleEmbedding, requirementEmbedding, skillsEmbedding)
%     \item Đảm bảo durability và backup
% \end{itemize}

% \vspace{0.3cm}
% \textbf{2. FAISS (In-Memory Cache):}
% \begin{itemize}
%     \item Load vectors từ PostgreSQL khi khởi động
%     \item Xây dựng IVF index với ncluster = sqrt(N)
%     \item Cập nhật realtime khi có CV/Job mới
%     \item Tốc độ truy vấn: 1-3ms
% \end{itemize}

% \end{frame}

%--- 2.4 LUỒNG TRUY VẤN ---
\begin{frame}{Cài đặt: Luồng truy vấn}
\begin{columns}[T]
\column{0.65\textwidth}
\justifying
\small

\textbf{Quy trình gợi ý việc làm cho CV gồm 5 bước:}

\vspace{0.3cm}
\begin{enumerate}
    \item \textbf{Load dữ liệu}: Tải danh sách CV và Job từ PostgreSQL với các vector nhúng đã tính trước
    
    \item \textbf{Kiểm tra thay đổi}: Kiểm tra mã băm nội dung (content hash) để xác định bản ghi thay đổi, chỉ tính lại embedding cho các bản ghi này
    
    \item \textbf{Xây dựng FAISS index}: Xây dựng hoặc tải chỉ mục FAISS từ các vector tiêu đề công việc
    
    \item \textbf{Áp dụng lọc phân tầng}: Chạy thuật toán lọc phân tầng 3 vòng để tìm $K_3$ công việc phù hợp nhất cho mỗi CV
    
    \item \textbf{Lưu kết quả}: Lưu kết quả gợi ý vào bảng RecommendJobForCV
\end{enumerate}

% \textbf{Input:} CV ID $\rightarrow$ Load vectors từ PostgreSQL

% \vspace{0.2cm}
% \textbf{Tham số:} $K_1=1000$, $K_2=300$, $K_3=30$ (cấu hình linh hoạt)

% \vspace{0.2cm}
% \textbf{Vòng 1 - Tiêu đề (FAISS):}
% \begin{itemize}
%     \item Tìm $K_1$ jobs có title tương đồng nhất
%     \item Sử dụng FAISS IVF index, độ tương đồng cosin
% \end{itemize}

% \vspace{0.15cm}
% \textbf{Vòng 2 - Kinh nghiệm:}
% \begin{itemize}
%     \item So sánh expEmb của CV với requirementEmb
%     \item Chọn $K_2$ jobs khớp nhất từ $K_1$ jobs
%     \item Skip nếu CV không có kinh nghiệm
% \end{itemize}

% \vspace{0.15cm}
% \textbf{Vòng 3 - Kỹ năng:}
% \begin{itemize}
%     \item So sánh skillsEmb với jobSkillsEmb
%     \item Trả về $K_3$ jobs cuối cùng
%     \item Skip nếu CV không có kỹ năng
% \end{itemize}

\column{0.3\textwidth}
\begin{figure}
    \centering
    \includegraphics[width=\textwidth]{figures/Truy vấn gợi ý.png}
    \caption{Luồng xử lý lọc phân tầng}
\end{figure}
\end{columns}
\end{frame}

%--- 2.5 TỐI ƯU HÓA ---
\begin{frame}{Cài đặt: Tối ưu hóa}
\begin{columns}[T]
\column{0.55\textwidth}
\justifying
\small

\textbf{Hệ thống áp dụng 3 chiến lược tối ưu chính:}

\vspace{0.1cm}
\textbf{1. Bộ nhớ đệm mã băm nội dung:}
\begin{itemize}
    \item Tính hash từ nội dung: MD5(JSON(title, skills, ...))
    \item Chỉ tính lại embedding khi hash thay đổi
    % \item Giảm 95\% tính toán nếu chỉ 5\% dữ liệu thay đổi
\end{itemize}

\vspace{0.1cm}
\textbf{2. Lọc công việc đang hoạt động:}
\begin{itemize}
    \item Chỉ xử lý jobs có status='ACTIVE'
    \item Loại bỏ jobs đã hết hạn (expiresAt < NOW)
    % \item Giảm khối lượng tính toán, đảm bảo kết quả hiệu lực
\end{itemize}

\vspace{0.1cm}
\textbf{3. Cập nhật định kỳ:}
\begin{itemize}
    \item Chu kỳ mặc định: 12 giờ (cấu hình linh hoạt)
    % \item Cân bằng độ tươi mới và tải hệ thống
\end{itemize}

\column{0.4\textwidth}
\begin{figure}
    \centering
    \includegraphics[width=\textwidth]{figures/Tối ưu hóa.png}
    \caption{Luồng xử lý bộ nhớ đệm mã băm nội dung}
\end{figure}
\end{columns}
\end{frame}

% PHẦN 3: ĐÁNH GIÁ

%--- 3.1 BỘ DỮ LIỆU ---
\begin{frame}{Đánh giá: Bộ dữ liệu}
\justifying

\textbf{Dataset thu thập:}
\begin{itemize}
    \item \textbf{Nguồn}: TopCV, VietnamWorks (crawled với permission)
    \item \textbf{Quy mô}: 5,000 CVs + 10,000 Jobs (tiếng Việt)
    \item \textbf{Phân chia}: 70\% train / 15\% validation / 15\% test
\end{itemize}

\vspace{0.3cm}
\textbf{Ground truth:}
\begin{itemize}
    \item Thu thập từ thực tế: CV đã ứng tuyển thành công vào Job
    \item Labeling thủ công: 3 HR experts đánh giá độ phù hợp (scale 1-5)
    \item Inter-rater agreement: Cohen's Kappa = 0.78 (substantial)
\end{itemize}

\vspace{0.3cm}
\textbf{Đặc điểm dữ liệu:}
\begin{itemize}
    \item Ngành nghề: IT (60\%), Marketing (20\%), Business (10\%), Others (10\%)
    \item Experience level: Intern-Junior (40\%), Mid (35\%), Senior+ (25\%)
\end{itemize}

\end{frame}

%--- 3.2 ĐỘ ĐO ---
\begin{frame}{Đánh giá: Độ đo}
\justifying

\textbf{1. Mean Reciprocal Rank (MRR):} Đo vị trí relevant item đầu tiên trong danh sách gợi ý.
$$MRR = \frac{1}{|Q|} \sum_{i=1}^{|Q|} \frac{1}{rank_i}$$

\vspace{0.3cm}
\textbf{2. Normalized Discounted Cumulative Gain (NDCG@K):} Đo chất lượng ranking, rel$_i$ là độ phù hợp tại vị trí i.
$$NDCG@K = \frac{DCG@K}{IDCG@K}, \quad DCG@K = \sum_{i=1}^{K} \frac{2^{rel_i} - 1}{\log_2(i+1)}$$

\vspace{0.3cm}
\textbf{3. Hit Rate@K:} Xác suất có ít nhất 1 kết quả phù hợp trong top-K.
$$HitRate@K = \frac{\text{Số queries có relevant item trong top-K}}{|Q|}$$


\end{frame}

%--- 3.3 KẾT QUẢ ---
\begin{frame}{Đánh giá: Kết quả thực nghiệm}
\justifying

\begin{table}[H]
\centering
\small
\begin{tabular}{|l|c|c|c|c|}
\hline
\textbf{Phương pháp} & \textbf{MRR} & \textbf{NDCG@10} & \textbf{Hit@10} & \textbf{Latency} \\
\hline
Keyword Matching & 0.42 & 0.51 & 0.68 & <10ms \\
\hline
TF-IDF + Cosine & 0.56 & 0.64 & 0.79 & ~20ms \\
\hline
mBERT + Flat & 0.71 & 0.76 & 0.89 & ~500ms \\
\hline
\textbf{PhoBERT-SimCSE + FAISS} & \textbf{0.78} & \textbf{0.82} & \textbf{0.93} & \textbf{~50ms} \\
\hline
+ Cascade (ours) & \textbf{0.81} & \textbf{0.85} & \textbf{0.95} & \textbf{~30ms} \\
\hline
\end{tabular}
\end{table}

\vspace{0.2cm}
\textbf{Nhận xét:}
\begin{itemize}
    \item PhoBERT vượt mBERT \textbf{+7\% MRR} nhờ hiểu tiếng Việt tốt hơn
    \item FAISS giảm latency từ 500ms xuống 50ms (10x)
    \item Cascade filtering tăng accuracy \textbf{+3\%} và giảm latency thêm 40\%
\end{itemize}

\end{frame}

% HẾT CHƯƠNG 3


% CHƯƠNG 4: KẾT QUẢ ĐẠT ĐƯỢC VÀ KẾT LUẬN
% File này chứa nội dung đã được sửa đổi theo yêu cầu

\makesection{Kết quả đạt được và Kết luận}

% 1. MÔI TRƯỜNG TRIỂN KHAI (BẢNG)

\begin{frame}{Công nghệ và Môi trường triển khai}
\justifying

\textbf{Kiến trúc triển khai:}

\begin{table}[H]
\centering
\small
\begin{tabular}{|p{4cm}|p{3.5cm}|p{3cm}|}
\hline
\textbf{Thành phần} & \textbf{Công nghệ} & \textbf{Nơi triển khai} \\
\hline
\textbf{Frontend} & React, TypeScript, Vite & Firebase Hosting \\
\hline
\textbf{Backend Server} & Node.js, Express & Google Cloud VM \\
\hline
\textbf{Database} & PostgreSQL, Prisma ORM & Google Cloud VM \\
\hline
\textbf{AI/Recommend Service} & Python, FastAPI, PhoBERT, FAISS & Google Cloud VM \\
\hline
\end{tabular}
\end{table}

\vspace{0.3cm}
\textbf{Dịch vụ ngoài:}
\begin{itemize}
    \item \textbf{Firebase Storage}: Lưu trữ file (CV, ảnh đại diện, logo công ty)
    \item \textbf{GitHub}: Quản lý mã nguồn
\end{itemize}

\end{frame}

% 2. TRIỂN KHAI CODE (CẤU TRÚC MÃ NGUỒN)

\begin{frame}[fragile]{Triển khai code - cấu trúc mã nguồn}
\justifying

\textbf{Kiến trúc:} Modular Clean Architecture, kết hợp ưu điểm
của Package by Feature và Layered Architecture.
\begin{columns}[T]
\column{0.5\textwidth}
\tiny
\vspace{0.2cm}
\textbf{Cấu trúc tổng thể:}
\begin{verbatim}
src/
├── app.ts                    # Cấu hình Express
├── server.ts                 # Điểm khởi động server
├── core/                     # Hạ tầng framework
│   ├── container/            # Dependency Injection
│   ├── middleware/           # Xác thực, xử lý lỗi
│   └── routes/               # Tổng hợp các route API
├── modules/                  # Các module chức năng
│   ├── user/                 # Quản lý người dùng
│   ├── company/              # Quản lý công ty
│   ├── job/                  # Quản lý tin tuyển dụng
│   ├── cv/                   # Quản lý CV
│   ├── application/          # Quản lý đơn ứng tuyển
│   └── notification/         # Quản lý thông báo
└── shared/                   # Phần dùng chung
    ├── constants/            # Hằng số, mã trạng thái
    ├── domain/               # Lỗi, interface dùng chung
    └── infrastructure/       # Kết nối database, dịch vụ ngoài
\end{verbatim}

\column{0.45\textwidth}
\vspace{0.2cm}
\tiny
\textbf{Mỗi module được tổ chức theo 4 tầng:}
\begin{verbatim}
{module}/
├── domain/                   # Tầng nghiệp vụ lõi
│   ├── entities/             # Thực thể miền
│   ├── enums/                # Kiểu liệt kê
│   └── repositories/         # Interface repository 
├── application/              # Tầng ứng dụng
│   ├── dtos/                 # Đối tượng truyền dữ liệu
│   ├── use-cases/            # Logic nghiệp vụ
│   └── helpers/              # Mapper, tiện ích
├── infrastructure/           # Tầng hạ tầng
│   ├── mappers/              # Chuyển đổi Entity <-> Model
│   └── repositories/         # Triển khai repository
└── interfaces/               # Tầng giao tiếp
    ├── controllers/          # Controller xử lý HTTP
    ├── routes/               # Định tuyến Express
    └── validators/           # Schema validate 
\end{verbatim}

% \column{0.45\textwidth}
% \begin{figure}
%     \centering
%     \includegraphics[width=\textwidth]{figures/structure.png}
%     \caption{4 tầng trong module User}
% \end{figure}
\end{columns}
\end{frame}

% 3. GIAO DIỆN 

\begin{frame}{Giao diện hệ thống - Ứng viên}
\begin{figure}[htb]
	\centering
	\includegraphics[width=0.8\linewidth]{figures/CandidateHome.png}
	\caption{Giao diện trang chủ của ứng viên}
\end{figure}

\end{frame}
\begin{frame}{Giao diện hệ thống - Ứng viên}

\begin{figure}[H]
	\centering
	\includegraphics[width=0.75\linewidth]{figures/RecommendJobs.png}
	\caption{Giao diện việc làm gợi ý}
\end{figure}
\end{frame}
\begin{frame}{Giao diện hệ thống - Ứng viên}
\begin{figure}[H]
	\centering
	\includegraphics[width=0.55\linewidth]{figures/mobile1.jpg}
	\caption{Giao diện trên điện thoại di động}
\end{figure}
\end{frame}
\begin{frame}{Giao diện hệ thống - Ứng viên}
\begin{figure}[H]
	\centering
	\includegraphics[width=0.7\linewidth]{figures/mobile2.jpg}
	\caption{Giao diện trên thiết bị máy tính bảng}
\end{figure}
\end{frame}

\begin{frame}{Giao diện hệ thống - Nhà tuyển dụng}
\begin{figure}
    \centering
    \includegraphics[width=0.8\textwidth]{figures/Recruiter Homepage.png}
    \caption{Dashboard nhà tuyển dụng quản lý đơn ứng tuyển}
\end{figure}
\end{frame}

\begin{frame}{Kết luận}
\justifying

\textbf{Các kết quả đạt được:}

\begin{itemize}
    \item Khảo sát các hệ thống tuyển dụng hiện có (VietnamWorks, TopCV, CareerViet)

    \item Phân tích yêu cầu hệ thống tuyển dụng trực tuyến

    \item Thiết kế hệ thống theo mô hình Clean Architecture với 6 module, 4 tầng

    \item Cài đặt hệ thống 3 thành phần (Backend, Frontend, AI Service) và triển khai trên cloud

    \item Xây dựng module gợi ý với PhoBERT, FAISS và thuật toán lọc ba tầng

    \item Đánh giá hệ thống với MRR, NDCG@10, Hit Rate
\end{itemize}

\end{frame}

\begin{frame}{Hạn chế và Hướng phát triển}
\justifying

\textbf{Hạn chế của đồ án:}

\begin{itemize}
    \item Mức cải thiện so với phương pháp cơ sở còn hạn chế (hai tầng lọc sau đóng góp 8\%)

    \item PhoBERT chưa được tinh chỉnh trên kho ngữ liệu tuyển dụng tiếng Việt

    \item Chưa triển khai quy trình CI/CD tự động

    \item Chưa đánh giá khả năng chịu tải và bảo mật trong môi trường thực tế
\end{itemize}

\vspace{0.3cm}
\textbf{Hướng phát triển:}

\begin{itemize}
    \item Tinh chỉnh PhoBERT và nghiên cứu learning-to-rank

    \item Triển khai CI/CD tự động (GitHub Actions/GitLab CI)

    \item Triển khai môi trường vận hành với giám sát, logging và kiểm toán bảo mật
\end{itemize}

\end{frame}

% Final PAGE
% Set the text that is showed on the final slide
\finalpagetext {Cảm ơn thầy cô đã chú ý lắng nghe!}
%----------------------------------------------------------------------------------------
\makefinalpage

%----------------------------------------------------------------------------------------
\end{document}
